\documentclass[10pt,a4paper]{amsart}
\usepackage{amsmath,amssymb,amsthm,graphicx}
% For commutative diagrams we use Xy-pic
\usepackage[all]{xy}

% We use 2cell for 2-commutative diagrams.
\xyoption{2cell}
% Theorem environments.
%
\theoremstyle{plain}
\newtheorem{theorem}[subsection]{Theorem}
\newtheorem{proposition}[subsection]{Proposition}
\newtheorem{lemma}[subsection]{Lemma}

\theoremstyle{definition}
\newtheorem{definition}[subsection]{Definition}
\newtheorem{example}[subsection]{Example}
\newtheorem{exercise}[subsection]{Exercise}
\newtheorem{situation}[subsection]{Situation}

\theoremstyle{remark}
\newtheorem{remark}[subsection]{Remark}
\newtheorem{remarks}[subsection]{Remarks}

\numberwithin{equation}{subsection}

% Macros
%
\def\lim{\mathop{\mathrm{lim}}\nolimits}
\def\colim{\mathop{\mathrm{colim}}\nolimits}
\def\Spec{\mathop{\mathrm{Spec}}}
\def\Hom{\mathop{\mathrm{Hom}}\nolimits}
\def\Ext{\mathop{\mathrm{Ext}}\nolimits}
\def\SheafHom{\mathop{\mathcal{H}\!\mathit{om}}\nolimits}
\def\SheafExt{\mathop{\mathcal{E}\!\mathit{xt}}\nolimits}
\def\Sch{\mathit{Sch}}
\def\Mor{\mathop{\mathrm{Mor}}\nolimits}
\def\Ob{\mathop{\mathrm{Ob}}\nolimits}
\def\Sh{\mathop{\mathit{Sh}}\nolimits}
\def\NL{\mathop{N\!L}\nolimits}
\def\CH{\mathop{\mathrm{CH}}\nolimits}
\def\proetale{{pro\text{-}\acute{e}tale}}
\def\etale{{\acute{e}tale}}
\def\QCoh{\mathit{QCoh}}
\def\Ker{\mathop{\mathrm{Ker}}}
\def\Im{\mathop{\mathrm{Im}}}
\def\Coker{\mathop{\mathrm{Coker}}}
\def\Coim{\mathop{\mathrm{Coim}}}

% Boxtimes
%
\DeclareMathSymbol{\boxtimes}{\mathbin}{AMSa}{"02}

%
% Macros for moduli stacks/spaces
%
\def\QCohstack{\mathcal{QC}\!\mathit{oh}}
\def\Cohstack{\mathcal{C}\!\mathit{oh}}
\def\Spacesstack{\mathcal{S}\!\mathit{paces}}
\def\Quotfunctor{\mathrm{Quot}}
\def\Hilbfunctor{\mathrm{Hilb}}
\def\Curvesstack{\mathcal{C}\!\mathit{urves}}
\def\Polarizedstack{\mathcal{P}\!\mathit{olarized}}
\def\Complexesstack{\mathcal{C}\!\mathit{omplexes}}
% \Pic is the operator that assigns to X its picard group, usage \Pic(X)
% \Picardstack_{X/B} denotes the Picard stack of X over B
% \Picardfunctor_{X/B} denotes the Picard functor of X over B
\def\Pic{\mathop{\mathrm{Pic}}\nolimits}
\def\Picardstack{\mathcal{P}\!\mathit{ic}}
\def\Picardfunctor{\mathrm{Pic}}
\def\Deformationcategory{\mathcal{D}\!\mathit{ef}}
\usepackage{cleveref}
\crefname{originallemma}{original lemma}{original lemmas}
\Crefname{originallemma}{Original lemma}{Original lemmas}
\makeatletter
\newlabel{lemma-higher-direct-images}{{03Q8}{source}{}{sourcepointer-1}{}}
\newlabel{lemma-higher-direct-images@cref}{{[originallemma][0][]03Q8}{[1][0][]source}{}{}{}}
\makeatother
\crefname{originallemma}{original lemma}{original lemmas}
\Crefname{originallemma}{Original lemma}{Original lemmas}
\makeatletter
\newlabel{lemma-describe-etale-local-ring}{{04HX}{source}{}{sourcepointer-2}{}}
\newlabel{lemma-describe-etale-local-ring@cref}{{[originallemma][0][]04HX}{[1][0][]source}{}{}{}}
\makeatother
\crefname{originallemma}{original lemma}{original lemmas}
\Crefname{originallemma}{Original lemma}{Original lemmas}
\makeatletter
\newlabel{limits-lemma-directed-inverse-system-affine-schemes-has-limit}{{01YW}{source}{}{sourcepointer-3}{}}
\newlabel{limits-lemma-directed-inverse-system-affine-schemes-has-limit@cref}{{[originallemma][0][]01YW}{[1][0][]source}{}{}{}}
\makeatother
\crefname{originallemma}{original lemma}{original lemmas}
\Crefname{originallemma}{Original lemma}{Original lemmas}
\makeatletter
\newlabel{lemma-directed-colimit-cohomology}{{03Q6}{source}{}{sourcepointer-4}{}}
\newlabel{lemma-directed-colimit-cohomology@cref}{{[originallemma][0][]03Q6}{[1][0][]source}{}{}{}}
\makeatother
\crefname{originallemma}{original lemma}{original lemmas}
\Crefname{originallemma}{Original lemma}{Original lemmas}
\makeatletter
\newlabel{lemma-colimit-complexes}{{0GIT}{source}{}{sourcepointer-5}{}}
\newlabel{lemma-colimit-complexes@cref}{{[originallemma][0][]0GIT}{[1][0][]source}{}{}{}}
\makeatother
\title{Stalks of higher direct images along a quasi-compact quasi-separated morphism as cohomology of the strict henselization}
\author{The Stacks Project Authors}
\date{stacks-project:etale-cohomology.tex:local-snapshot-20261011; GNU FDL 1.2 or later}
\begin{document}
\maketitle
\noindent\textit{Editorial note.} All delivered mathematical text is the author's own, copied from the pinned Stacks Project chapter file etale-cohomology.tex; no statement or proof step was rewritten, reordered or removed.\par\medskip
\noindent\textit{Editorial note.} The proof is the author's complete argument for the sheaf statement and for the derived statement, the latter obtained from the former by choosing a representative of the bounded below complex.\par\medskip
\noindent\textit{Editorial note.} Every labelled reference occurring inside the counted spans is carried in assembly.external\_references as a non-load-bearing navigation pointer (same-chapter ones at their definition site, cross-chapter ones at the invocation inside the counted span), so no load-bearing mathematics is replaced by a stub. No bibliography citation occurs inside the counted spans.\par\medskip
\noindent\textit{Editorial note.} The checked chapter list for same-name collisions is the set of stacks chapters staged in this repository: chow, cohomology, divisors, etale-cohomology, flat, more-algebra, more-morphisms, morphisms, spaces-more-morphisms, varieties together with this card's chapter etale-cohomology.\par\medskip
\noindent\textit{Editorial note.} The preamble is the chapter's own native preamble excerpt already used for this source: the Xy-pic setup (native preamble lines 14-18) and the native theorem declarations with the mathematical macros (native preamble lines 154-217).\par\medskip
\noindent\textit{Allowed standard prerequisites (editorial).}\begin{itemize}
\item Strictly henselian local ring at a geometric point and the small etale site of its spectrum. Same chapter, the strictly henselian local ring O\textasciicircum{}\{sh\}\_\{S,sbar\} at a geometric point, its spectrum as a limit of etale neighbourhoods, and cohomology over it. Supplies the right hand side of the displayed identification.
\item Cohomology of a limit of schemes and the comparison with the limit of cohomologies. Same chapter, the result that etale cohomology commutes with the filtered limit defining the strict henselization. Supplies the passage from the colimit of the cohomologies over etale neighbourhoods to cohomology over the strict henselization.
\item Stalks of higher direct images computed by the Cech-to-cohomology spectral sequence. Same chapter, the description of the stalk of R\textasciicircum{}n f\_*F at a geometric point by the Cech cohomology of the covering given by a cofinal system of etale neighbourhoods. Supplies the left hand side and the spectral sequence with which the proof compares the two sides.
\end{itemize}\medskip
\ifdefined\hypertarget\hypertarget{sourcepointer-1}{}\fi
\noindent\textit{Source reference.} Original lemma 03Q8: sources/stacks-etale-cohomology/pinned-original-source.tex, source lines 7510--7510. Navigation pointer invoked at proof line 8042; the counted span contains the pointer only and the cited result is not reproduced inside it, so the pointer is registered as a non-load-bearing navigation stub in the established form for Stacks material. No bibliography citation is involved.\par\medskip
\ifdefined\hypertarget\hypertarget{sourcepointer-2}{}\fi
\noindent\textit{Source reference.} Original lemma 04HX: sources/stacks-etale-cohomology/pinned-original-source.tex, source lines 4257--4257. Navigation pointer invoked at proof line 8055; the counted span contains the pointer only and the cited result is not reproduced inside it, so the pointer is registered as a non-load-bearing navigation stub in the established form for Stacks material. No bibliography citation is involved.\par\medskip
\ifdefined\hypertarget\hypertarget{sourcepointer-3}{}\fi
\noindent\textit{Source reference.} Original lemma 01YW: sources/stacks-etale-cohomology/pinned-original-source.tex, source lines 8057--8057. Navigation pointer invoked at proof line 8057; the counted span contains the pointer only and the cited result is not reproduced inside it, so the pointer is registered as a non-load-bearing navigation stub in the established form for Stacks material. No bibliography citation is involved.\par\medskip
\ifdefined\hypertarget\hypertarget{sourcepointer-4}{}\fi
\noindent\textit{Source reference.} Original lemma 03Q6: sources/stacks-etale-cohomology/pinned-original-source.tex, source lines 7464--7464. Navigation pointer invoked at proof line 8063; the counted span contains the pointer only and the cited result is not reproduced inside it, so the pointer is registered as a non-load-bearing navigation stub in the established form for Stacks material. No bibliography citation is involved.\par\medskip
\ifdefined\hypertarget\hypertarget{sourcepointer-5}{}\fi
\noindent\textit{Source reference.} Original lemma 0GIT: sources/stacks-etale-cohomology/pinned-original-source.tex, source lines 7853--7853. Navigation pointer invoked at proof line 8065; the counted span contains the pointer only and the cited result is not reproduced inside it, so the pointer is registered as a non-load-bearing navigation stub in the established form for Stacks material. No bibliography citation is involved.\par\medskip
\begin{theorem}
\label{theorem-higher-direct-images}
Let $f: X \to S$ be a quasi-compact and quasi-separated morphism of schemes,
$\mathcal{F}$ an abelian sheaf on $X_\etale$, and $\overline{s}$ a
geometric point of $S$ lying over $s \in S$. Then
$$
\left(R^nf_* \mathcal{F}\right)_{\overline{s}} =
H_\etale^n( X \times_S \Spec(\mathcal{O}_{S, \overline{s}}^{sh}),
p^{-1}\mathcal{F})
$$
where $p : X \times_S \Spec(\mathcal{O}_{S, \overline{s}}^{sh}) \to X$
is the projection. For $K \in D^+(X_\etale)$ and $n \in \mathbf{Z}$
we have
$$
\left(R^nf_*K\right)_{\overline{s}} =
H_\etale^n(X \times_S \Spec(\mathcal{O}_{S, \overline{s}}^{sh}), p^{-1}K)
$$
In fact, we have
$$
\left(Rf_*K\right)_{\overline{s}}
=
R\Gamma_\etale(X \times_S \Spec(\mathcal{O}_{S, \overline{s}}^{sh}), p^{-1}K)
$$
in $D^+(\textit{Ab})$.
\end{theorem}

\begin{proof}
Let $\mathcal{I}$ be the category of \'etale neighborhoods of $\overline{s}$
on $S$. By Lemma \ref{lemma-higher-direct-images}
we have
$$
(R^nf_*\mathcal{F})_{\overline{s}} =
\colim_{(V, \overline{v}) \in \mathcal{I}^{opp}}
H_\etale^n(X \times_S V, \mathcal{F}|_{X \times_S V}).
$$
We may replace $\mathcal{I}$ by the initial subcategory consisting
of affine \'etale neighbourhoods of $\overline{s}$. Observe that
$$
\Spec(\mathcal{O}_{S, \overline{s}}^{sh}) =
\lim_{(V, \overline{v}) \in \mathcal{I}} V
$$
by Lemma \ref{lemma-describe-etale-local-ring} and
Limits, Lemma
\ref{limits-lemma-directed-inverse-system-affine-schemes-has-limit}.
Since fibre products commute with limits we also obtain
$$
X \times_S \Spec(\mathcal{O}_{S, \overline{s}}^{sh}) =
\lim_{(V, \overline{v}) \in \mathcal{I}} X \times_S V
$$
We conclude by Lemma \ref{lemma-directed-colimit-cohomology}.
For the second variant, use the same argument using
Lemma \ref{lemma-colimit-complexes} instead of
Lemma \ref{lemma-directed-colimit-cohomology}.

\medskip\noindent
To see that the last statement is true, it suffices to produce a map
$\left(Rf_*K\right)_{\overline{s}} \to
R\Gamma_\etale(X \times_S \Spec(\mathcal{O}_{S, \overline{s}}^{sh}), p^{-1}K)$
in $D^+(\textit{Ab})$ which realizes the ismorphisms on cohomology
groups in degree $n$ above for all $n$. To do this, choose a
bounded below complex $\mathcal{J}^\bullet$ of injective abelian
sheaves on $X_\etale$ representing $K$. The complex
$f_*\mathcal{J}^\bullet$ represents $Rf_*K$. Thus the complex
$$
(f_*\mathcal{J}^\bullet)_{\overline{s}} =
\colim_{(V, \overline{v}) \in \mathcal{I}^{opp}}
(f_*\mathcal{J}^\bullet)(V)
$$
represents $(Rf_*K)_{\overline{s}}$. For each $V$ we have maps
$$
(f_*\mathcal{J}^\bullet)(V) =
\Gamma(X \times_S V, \mathcal{J}^\bullet)
\longrightarrow
\Gamma(X \times_S \Spec(\mathcal{O}_{S, \overline{s}}^{sh}),
p^{-1}\mathcal{J}^\bullet)
$$
and the target complex represents
$R\Gamma_\etale(X \times_S \Spec(\mathcal{O}_{S, \overline{s}}^{sh}), p^{-1}K)$
in $D^+(\textit{Ab})$.
Taking the colimit of these maps we obtain the result.
\end{proof}

\end{document}
